\section{The degree-two obstruction over arbitrary characteristic-zero fields}
\label{app:degree-two}

This appendix supplies the scalar-extension and descent details used in
Lemma~\ref{lem:degree-two}.  Let
\[
 F=(F_1,\ldots,F_m)\in K[X_1,\ldots,X_m]^m
 \quad\text{with}\quad
 \det\left(\frac{\partial F_i}{\partial X_j}\right)\in K^*,
\]
and suppose for contradiction that its function-field degree is two.

\paragraph{The external input.}
Campbell's unnumbered theorem
\cite[Theorem, p.~244]{Campbell} states that a polynomial map
\(H:\mathbb C^m\to\mathbb C^m\) has a polynomial inverse if its Jacobian
determinant never vanishes and the induced extension
\[
 \mathbb C(H_1,\ldots,H_m)\subset
 \mathbb C(X_1,\ldots,X_m)
\]
is normal.  In characteristic zero, normality here is equivalent to the
extension being Galois.  Razar and Wright later gave algebraic treatments of
this Galois case \cite{Razar,Wright}; the argument below needs only
Campbell's complex theorem.

\paragraph{Generic degree survives scalar extension.}
Let \(K_0\subset K\) be the subfield generated over \(\mathbb Q\) by the
finitely many coefficients of the \(F_i\), and write \(F_0\) for the same
formulas over \(K_0\).  The nonzero Jacobian determinant makes
\begin{equation}
 \varphi_0:K_0[Y_1,\ldots,Y_m]\longrightarrow
 K_0[X_1,\ldots,X_m],
 \qquad Y_i\longmapsto F_i,
 \label{eq:degree-two-coordinate-map}
\end{equation}
injective and generically finite.  Indeed, the Jacobian criterion in
characteristic zero makes the \(F_i\) algebraically independent, so both
fraction fields have transcendence degree \(m\).  If their extension has
degree \(d_0\), there is a nonempty principal open \(U_0\) of the target on
which \(F_0\) is finite locally free of rank \(d_0\): clear denominators in
the finite extension and apply generic freeness.  For every field extension
\(E/K_0\), base change over \((U_0)_E\) is again finite locally free of rank
\(d_0\).  Taking generic fibers gives
\begin{equation}
 [E(X_1,\ldots,X_m):E(F_1,\ldots,F_m)]=d_0.
 \label{eq:degree-two-base-change}
\end{equation}
Taking \(E=K\) shows that \(d_0=2\).

\paragraph{The quadratic extension is Galois.}
Because \(K_0\) is finitely generated over \(\mathbb Q\), choose an embedding
\(K_0\hookrightarrow\mathbb C\).  Equation
\eqref{eq:degree-two-base-change} with \(E=\mathbb C\) says that
\[
 \mathbb C(X_1,\ldots,X_m)/
 \mathbb C(F_1,\ldots,F_m)
\]
has degree two.  It is separable, and every separable quadratic field
extension is normal: the second root of the minimal polynomial of an element
is its trace minus that element.  The complexified map therefore satisfies
Campbell's hypotheses and is a polynomial automorphism.

\paragraph{Descent of the inverse.}
The base change of \eqref{eq:degree-two-coordinate-map} along
\(K_0\hookrightarrow\mathbb C\) is an isomorphism.  Since \(\mathbb C\) is
faithfully flat over \(K_0\), tensoring the kernel and cokernel of
\(\varphi_0\) with \(\mathbb C\) detects their vanishing.  Hence
\(\varphi_0\) is an isomorphism.  After base change to \(K\), the original
\(F\) is a polynomial automorphism, so its function-field degree is one,
contrary to the assumption.
